




\begin{definition}[Functor of points]\label{nhq-functor-of-points}\source{nitsure-hilbert-quot:page-0001}
Every scheme $X$ defines a contravariant functor $h_X$, called its functor
of points, from schemes to sets by $h_X(T)=Mor(T,X)$. The scheme $X$ is
determined up to unique isomorphism by $h_X$ via the Yoneda lemma; its
restriction to affine schemes already suffices.
\end{definition}

\begin{definition}[Fpqc descent for set-valued functors]\label{nhq-fpqc-descent}\source{nitsure-hilbert-quot:page-0001}
For a set-valued functor $F$, descent in the fpqc topology is the exactness
of
$$F(U)\longrightarrow \prod_i F(U_i)
\stackrel{\to}{\scriptstyle\to}\prod_{i,j}F(U_i\times_U U_j)$$
for every fpqc open cover $(U_i\to U)$. Every representable functor
satisfies this descent condition, so it is necessary for representability,
though it is not sufficient.
\end{definition}

\centerline{\large\bf The Functors $\Hilb_{\PP^n}$}


\medskip

\begin{definition}[Hilbert functor]\label{nhq-hilbert-functor}\source{nitsure-hilbert-quot:page-0002}\uses{nhq-functor-of-points}
The main problem addressed in this series of lectures, in its simplest
form, is as follows. If $S$ is a locally noetherian scheme,
a \textbf{family of subschemes of $\PP^n$ parametrised by $S$} will mean
a closed subscheme $Y \subset \PP^n_S = \PP^n_{\Z}\times S$
such that $Y$ is flat over $S$.
If $f: T\to S$ is any morphism of locally noetherian schemes, then
by pull-back we get a family $f^*(Y) = (\id \times f)^{-1}(Y)
\subset \PP^n_T$
parametrised by $T$, from a family $Y$ parametrised by $S$.
This defines a contravariant functor $\Hilb_{\PP^n}$
from the category of all locally noetherian schemes to the category of sets,
which associates to any $S$ the set of all such families
$$\Hilb_{\PP^n}(S) = \{ Y \subset \PP^n_S \,|\,
Y \mbox{ is flat over }S \}$$
\end{definition}

\textbf{Question:} Is the functor $\Hilb_{\PP^n}$ representable?

\begin{theorem*}\label{nhq-hilbert-representability}\source{nitsure-hilbert-quot:page-0002}\uses{nhq-hilbert-functor}
Grothendieck proved that this question has an affirmative answer, that is,
there exists a locally noetherian scheme $\hilb_{\PP^n}$ together with
a family $Z \subset \PP^n_{\Z}\times \hilb_{\PP^n}$
parametrised by $\hilb_{\PP^n}$,
such that any family $Y$ over $S$ is obtained as the
pull-back of $Z$
by a uniquely determined morphism $\varphi_Y : S \to \hilb_{\PP^n}$.
In other words, $\Hilb_{\PP^n}$ is isomorphic to
the functor $Mor(-,\hilb_{\PP^n})$.
\end{theorem*}

\begin{proof}
\notready The source states Grothendieck's representability theorem here but does not supply a proof.
\end{proof}





\medskip


\bigskip


\centerline{\large\bf The Functors $\Quot_{\oplus^r {\mathcal O}_{\PP^n}}$}
\nopagebreak
\medskip
\nopagebreak
A family $Y$ of subschemes of $\PP^n$ parametrised by $S$
is the same as
a coherent quotient sheaf $q: {\mathcal O}_{\PP^n_S} \to {\mathcal O}_Y$ on
$\PP^n_S$, such that ${\mathcal O}_Y$ is flat over $S$.
This way of looking at the functor $\Hilb_{\PP^n}$
has the following fruitful generalisation.


\begin{definition}[Quotient family on projective space]\label{nhq-projective-quot-family}\source{nitsure-hilbert-quot:page-0003}\uses{nhq-hilbert-functor}
Let $r$ be any positive integer.
A \textbf{family of quotients of $\oplus^r {\mathcal O}_{\PP^n}$
parametrised by a locally noetherian scheme $S$} will mean
a pair $(\F,q)$ consisting of

(i) a coherent sheaf $\F$ on $\PP^n_S$ which is flat over $S$, and

(ii) a surjective ${\mathcal O}_{\PP^n_S}$-linear homomorphism of sheaves
$q:\oplus^r {\mathcal O}_{\PP^n_S}  \to \F$.

Two such families $(\F,q)$ and $(\F,q)$ parametrised by $S$ will
be regarded as equivalent if there exists an isomorphism $f: \F \to \F'$
which takes $q$ to $q'$, that is, the following diagram commutes.
$$\begin{array}{ccc}
\oplus^r {\mathcal O}_{\PP^n} &\stackrel{q}{\to} & \F \\
\|                  &     & ~ \dna {\scriptstyle f} \\
\oplus^r {\mathcal O}_{\PP^n} &\stackrel{q'}{\to} & \F'
\end{array}$$
This is the same as the condition
$\ker(q) = \ker(q')$. We will denote by $\langle \F, q \rangle$
an equivalence class. If $f : T\to S$ is a morphism of
locally noetherian schemes, then pulling back the quotient
$q: \oplus^r{\mathcal O}_{\PP^n_S} \to \F$
under $\id\times f : \PP^n_T \to \PP^n_S$
defines a family $f^*(q) : \oplus^r {\mathcal O}_{\PP^n_T} \to f^*(\F)$ over $T$,
which makes sense as tensor product is right-exact and preserves flatness.
The operation of pulling back respects equivalence of families,
therefore it gives rise to a contravariant functor
$\Quot_{\oplus^r {\mathcal O}_{\PP^n}}$ from the category of
all locally noetherian schemes to the category of sets,
by putting
$$\Quot_{\oplus^r {\mathcal O}_{\PP^n}}(S) =
\{\mbox{ All } \langle \F, q \rangle \mbox{ parametrised by }S\}$$
\end{definition}


\begin{theorem*}[Representability of the projective Quot functor]\label{nhq-projective-quot-representability}\source{nitsure-hilbert-quot:page-0003}\uses{nhq-projective-quot-family}
It is immediate that the functor $\Quot_{\oplus^r {\mathcal O}_{\PP^n}}$
satisfies faithfully flat descent. It was proved by
Grothendieck that in fact the above functor
is representable on the category of all locally noetherian schemes
by a scheme $\quot_{\oplus^r{\mathcal O}_{\PP^n}}$.
\end{theorem*}

\begin{proof}
\notready The source records Grothendieck's theorem without a proof at this point.
\end{proof}


\medskip


\bigskip

\centerline{\large\bf The Functors $\Hilb_{X/S}$ and $\Quot_{E/X/S}$}

\medskip

\begin{definition}[Relative Quot and Hilbert functors]\label{nhq-relative-quot-family}\source{nitsure-hilbert-quot:page-0004}\uses{nhq-projective-quot-family}
The above functors $\Hilb_{\PP^n}$ and $\Quot_{\oplus^r {\mathcal O}_{\PP^n}}$
admit the following simple generalisations.
Let $S$ be a noetherian scheme and let $X\to S$ be a finite type
scheme over it. Let $E$ be a coherent sheaf on $X$.
Let $Sch_S$ denote the category of all locally noetherian schemes
over $S$. For any $T\to S$ in $Sch_S$, a \textbf{family of quotients of $E$
parametrised by $T$} will mean
a pair $(\F,q)$ consisting of

(i) a coherent sheaf $\F$ on $X_T = X\times_ST$
such that the schematic support of $\F$ is proper over $T$
and $\F$ is flat over $T$, together with

(ii) a surjective ${\mathcal O}_{X_T}$-linear homomorphism of sheaves
$q:E_T  \to \F$ where $E_T$ is the pull-back of $E$ under
the projection $X_T \to X$.

Two such families $(\F,q)$ and $(\F,q)$ parametrised by $T$ will
be regarded as equivalent if $\ker(q) = \ker(q')$),
and $\langle \F, q \rangle$ will denote
an equivalence class. Then as properness and flatness are preserved
by base-change, and as tensor-product is right exact,
the pull-back of $\langle \F, q \rangle$ under an $S$-morphism $T'\to T$ is
well-defined, which gives
a set-valued contravariant functor
$\Quot_{E/X/S} : Sch_S \to Sets$ under which
$$T\mapsto  \{\mbox{ All } \langle \F, q \rangle \mbox{ parametrised by }T\}$$

When $E = {\mathcal O}_X$, the functor $\Quot_{{\mathcal O}_X/X/S} : Sch_S \to Sets$
associates to $T$ the set of all closed subschemes $Y\subset X_T$
that are proper and flat over $T$. We denote this functor by $\Hilb_{X/S}$.

Note in particular that we have
$$\Hilb_{\PP^n} = \Hilb_{\PP^n_{\Z}/\spec \Z} \mbox{ and }
\Quot_{\oplus^r {\mathcal O}_{\PP^n}} =
\Quot_{\oplus^r {\mathcal O}_{\PP^n_{\Z}}/\PP^n_{\Z}/\spec \Z}$$


It is clear that the functors $\Quot_{E/X/S}$ and $\Hilb_{X/S}$
satisfy faithfully flat descent,
so it makes sense to pose the question of their
representability.
\end{definition}



\medskip


\bigskip

\centerline{\large\bf Stratification by Hilbert Polynomials}

\medskip


\begin{definition}[Hilbert polynomial and stratification]\label{nhq-hilbert-polynomial}\source{nitsure-hilbert-quot:page-0005}\uses{nhq-relative-quot-family}
Let $X$ be a finite type scheme over
a field $k$, together with a line bundle $L$.
Recall that if $F$ is a coherent sheaf on
$X$ whose support is proper over $k$, then
the \textbf{Hilbert polynomial} ${\Phi}\in \Q[\lambda]$ of $F$ is defined by
the function
$${\Phi}(m) = \chi(F(m)) =
\sum_{i=0}^n (-1)^i \dim_k H^i(X, F\otimes L^{\otimes m})$$
where the dimensions of the cohomologies are finite because of the
coherence and properness conditions. The fact that $\chi(F(m))$
is indeed a polynomial in $m$ under the above assumption is a special
case of what is known as Snapper's Lemma (see Kleiman [K] for a proof).



Let $X\to S$ be a finite type morphism of noetherian schemes,
and let $L$ be a line bundle on $X$. Let $\F$
be any coherent sheaf on
$X$ whose schematic support is proper over $S$. Then for each $s\in S$,
we get a polynomial ${\Phi}_s \in  \Q[\lambda]$
which is the Hilbert polynomial of
the restriction $\F_s = F|_{X_s}$ of $\F$ to the fiber $X_s$ over $s$,
calculated with respect to the line bundle $L_s = L|_{X_s}$.
If $\F$ is flat over $S$ then the function $s\mapsto {\Phi}_s$
from the set of points of $S$ to the polynomial ring $\Q[\lambda]$ is
known to be locally constant on $S$.


This shows that the functor $\Quot_{E/X/S}$ naturally decomposes
as a co-product
$$\Quot_{E/X/S} = \coprod_{{\Phi}\in \Q[\lambda]}\, \Quot_{E/X/S}^{\Phi,L}$$
where for any polynomial ${\Phi}\in \Q[\lambda]$, the functor
$\Quot_{E/X/S}^{\Phi,L}$ associates
to any $T$ the set of all equivalence classes of
families $\langle \F,q \rangle$ such that at each $t\in T$ the
Hilbert polynomial of the restriction $\F_t$, calculated
using the pull-back of $L$, is ${\Phi}$. Correspondingly,
the representing scheme  $\quot_{E/X/S}$, when it exists, naturally decomposes
as a co-product
$$\quot_{E/X/S} = \coprod_{{\Phi}\in \Q[\lambda]}\, \quot_{E/X/S}^{\Phi,L}$$

\textbf{Note }
We will generally take $X$ to be (quasi-)projective over $S$,
and $L$ to be a relatively very ample line bundle. Then indeed
the Hilbert and Quot functors are representable by schemes,
but not in general.
\end{definition}



\medskip


\bigskip

\centerline{\large\bf Construction of Grassmannian}


\medskip

The following explicit construction of
the Grassmannian scheme $\grass(r,d)$ over $\Z$ is best understood
as the construction of a quotient
$GL_{d,\Z}\backslash V$, where $V$ is the scheme of all
$d\times r$ matrices of rank $d$, and the group-scheme
$GL_{d,\Z}$ acts on $V$ on the left by matrix multiplication. However,
we will not use the language of group-scheme actions
here, instead, we give a direct elementary construction of
the Grassmannian scheme.

The reader can take $d=1$ in what follows, in a first reading, to
get the special case $\grass(r,1) = \PP^{r-1}_{\Z}$, which
has another construction as $\proj \Z[x_1,\ldots,x_r]$.



\medskip



\begin{construction}[Grassmannian by affine gluing]\label{nhq-grassmannian-construction}\source{nitsure-hilbert-quot:page-0007}
\textbf{Construction by gluing together affine patches }
For any integers $r\ge d\ge 1$, the Grassmannian scheme
$\grass(r,d)$ over $\Z$, together with the tautological
quotient $u:\oplus^r\,{\mathcal O}_{\grass(r,d)} \to \U$ where \
$\U$ is a rank $d$ locally free sheaf on $\grass(r,d)$,
can be explicitly constructed as follows.



If $M$ is a $d\times r$-matrix, and
$I\subset \{1,\ldots, r\}$ with cardinality $\#(I)$ equal
to $d$, the $I$ th minor $M_I$ of $M$ will
mean the $d\times d$ minor of $M$ whose columns are indexed by
$I$.


For any subset $I\subset \{1,\ldots, r\}$ with $\#(I) =d$,
consider the $d\times r$ matrix $X^I$ whose $I$ the minor
$X^I_I$ is the $d\times d$ identity matrix $1_{d\times d}$, while the remaining
entries of $X^I$ are independent variables $x^I_{p,q}$ over $\Z$.
Let $\Z[X^I]$ denote the polynomial ring in the variables
$x^I_{p,q}$, and let $U^I = \spec \Z[X^I]$, which is non-canonically
isomorphic to the affine space $\AA^{d(r-d)}_{\Z}$.

For any $J\subset \{1,\ldots, r\}$ with $\#(J) =d$, let
$P^I_J = \det(X^I_J) \in \Z[X^I]$
where $X^I_J$ is the $J$ th minor of $X^I$.
Let $U^I_J = \spec \Z[X^I, 1/P^I_J]$
the open subscheme of $U^I$ where $P^I_J$ is invertible. This means
the $d\times d$-matrix $X^I_J$ admits an inverse $(X^I_J)^{-1}$
on $U^I_J$.

For any $I$ and $J$, a ring homomorphism
$\theta_{I,J} : \Z[X^J, 1/P^J_I] \to \Z[X^I, 1/P^I_J]$ is
defined as follows. The images of the variables
$x^J_{p,q}$ are given by the entries of the matrix formula
$\theta_{I,J}(X^J) = (X^I_J) ^{-1}\,X^I$.
In particular, we have $\theta_{I,J}(P^J_I) = 1/P^I_J$,
so the map extends to $\Z[X^J, 1/P^J_I]$.

Note that $\theta_{I,I}$ is identity on $U^I_I = U^I$, and
we leave it to the reader to verify that \notready
for any three subsets $I$, $J$ and $K$ of $\{ 1,\ldots, r\}$ of
cardinality $d$, the co-cycle condition
$\theta_{I,K}= \theta_{I,J}\theta_{J,K}$
is satisfied. Therefore the schemes
$U^I$, as $I$ varies over all the
${r\choose d}$ different
subsets of  $\{ 1,\ldots, r\}$ of cardinality $d$,
can be glued together by the co-cycle $(\theta_{I,J})$ to form a
finite-type scheme $\grass(r,d)$ over $\Z$. As each $U^I$ is
isomorphic to $\AA^{d(r-d)}_{\Z}$, it follows that
$\grass(r,d) \to \spec \Z$ is smooth of relative dimension $d(r-d)$.



\medskip

\textbf{Separatedness }
The intersection of the diagonal of $\grass(r,d)$ with $U^I\times U^J$
can be seen to be the closed subscheme $\Delta_{I,J}\subset U^I\times U^J$
defined by entries of the matrix formula $X^J_I\,X^I - X^J =0$,
and so $\grass(r,d)$ is a separated scheme.


\medskip

\textbf{Properness }
We now show that $\pi : \grass(r,d) \to \spec \Z$ is proper. It is enough to
verify the valuative criterion of properness for discrete valuation rings.
Let $\Rr$ be a dvr, $\K$ its quotient field, and let
$\varphi : \spec \K \to \grass(r,d)$ be a morphism. This is given by
a ring homomorphism $f : \Z[X^I] \to \K$ for some $I$. Having fixed
one such $I$, next choose $J$ such that $\nu(f(P^I_J))$ is minimum,
where $\nu : \K \to \Z\bigcup \{ \infty\}$ denotes
the discrete valuation. As $P^I_I = 1$, note that
$\nu(f(P^I_J)) \le 0$, therefore
$f(P^I_J) \ne 0$ in $\K$ and so the matrix $f(X^I_J)$
lies in $GL_d(\K)$.

Now consider the homomorphism
$g : \Z[X^J] \to \K$ defined by entries of the matrix formula
$$g(X^J) = f((X^I_J)^{-1}\,X^I)$$
Then $g$ defines the same morphism $\varphi : \spec \K \to \grass(r,d)$,
and moreover all $d\times d$ minors $X^J_K$ satisfy
$\nu(g(P^J_K)) \ge 0$.
As the minor $X^J_J$ is identity, it follows from the above that in fact
$\nu(g(x^J_{p,q}))\ge 0$ for all entries of $X^J$.
Therefore, the map $g : \Z[X^J] \to \K$
factors uniquely via $\Rr \subset \K$. The resulting morphism of
schemes $\spec \Rr \to U^J \hra \grass(r,d)$ prolongs
$\varphi : \spec \K \to \grass(r,d)$. We have already checked
separatedness of $\grass(r,d)$, so now we see that
$\grass(r,d) \to \spec \Z$ is proper.


\medskip



\textbf{Universal quotient }
We next define a rank $d$ locally free sheaf $\U$ on $\grass(r,d)$
together with a surjective homomorphism $\oplus^r\,{\mathcal O}_{\grass(r,d)} \to \U$.
On each $U^I$ we define a surjective homomorphism
$u^I :\oplus^r\,{\mathcal O}_{U^I} \to  \oplus^d\,{\mathcal O}_{U^I}$
by the matrix $X^I$. Compatible with the co-cycle $(\theta_{I,J})$
for gluing the affine pieces $U^I$, we give gluing data
$(g_{I,J})$ for gluing together the trivial bundles $\oplus^d\,{\mathcal O}_{U^I}$
by putting
$$g_{I,J} = (X^I_J)^{-1} \in GL_d(U^I_J)$$
This is compatible with the homomorphisms $u^I$,
so we get a surjective homomorphism
$u : \oplus^r\,{\mathcal O}_{\grass(r,d)} \to \U$.


\medskip
\textbf{Projective embedding }
As $\U$ is given by the transition functions $g_{I,J}$ described above,
the determinant line bundle
$\det(\U)$ is given by the transition functions
$\det(g_{I,J}) = 1/P^I_J \in GL_1(U^I_J)$.
For each $I$, we define a global section
$$\sigma_I \in \Gamma(\grass(r,d),\det(\U))$$
by putting $\sigma_I|_{U^J} = P^J_I \in
\Gamma(U^J,{\mathcal O}_{U^J})$ in terms of the trivialization over the open cover
$(U^J)$.
We leave it to the reader to verify that the sections $\sigma_I$ form \notready
a linear system which is base point free and separates points
relative to $\spec \Z$, and so gives an embedding of
$\grass(r,d)$ into $\PP^m_{\Z}$ where $m = {r\choose d}-1$.
This is a closed embedding by the properness of
$\pi:\grass(r,d) \to \spec \Z$. In particular, $\det(\U)$
is a relatively very ample line bundle on $\grass(r,d)$ over $\Z$.
\end{construction}


\medskip

{\footnotesize
\textbf{Note } The $\sigma_I$ are known as the Pl\"ucker coordinates,
and these satisfy certain quadratic polynomials known as the Pl\"ucker
relations, which define the projective image of the Grassmannian.
We will not need these facts. }%end of footnotesize
